% SPDX-License-Identifier: CC-BY-4.0
\documentclass[11pt,letterpaper]{article}

\usepackage[T1]{fontenc}
\usepackage{amsmath,amssymb}
\usepackage{tikz}
\usepackage{pgfplots}
\pgfplotsset{compat=1.18}
\usepackage{algorithm}
\usepackage{algpseudocode}
\usepackage{booktabs}
\usepackage{enumitem}
\usepackage[hidelinks]{hyperref}

\newcommand{\N}{\mathbb{N}}
\newcommand{\Z}{\mathbb{Z}}
\newcommand{\Q}{\mathbb{Q}}
\newcommand{\R}{\mathbb{R}}
\newcommand{\C}{\mathbb{C}}
\newcommand{\Map}[3]{#1 \colon #2 \to #3}
\newcommand{\Mod}[1]{#1\text{-}\mathrm{Mod}}
\newcommand{\GroupHom}{\mathrm{Hom}_{\mathbf{Grp}}}
\newcommand{\RingHom}{\mathrm{Hom}_{\mathbf{Ring}}}

\newtheorem{theorem}{Theorem}
\newtheorem{lemma}[theorem]{Lemma}
\newtheorem{corollary}[theorem]{Corollary}
\newtheorem{definition}[theorem]{Definition}
\newtheorem{remark}[theorem]{Remark}
\newenvironment{proof}{\noindent\textit{Proof.}}{\hfill\(\square\)}

\title{Generalized Collatz Library}
\author{ascendnoosphere}
\date{\today}

\begin{document}

\maketitle

\noindent
This document studies the raw map
\[
T_{q,D}(n)=\begin{cases} n/D, & D\mid n,\\ qn+1, & D\nmid n,\end{cases}
\]
on the slice whose divisor is the repunit sum
\(S_k=S_k(q)=1+q+\cdots+q^{k-1}\). The offset is one, so the affine branch is
\(n\mapsto qn+1\), and the divisor identity \((q-1)S_k=q^k-1\) exhibits
\(S_k\) as a divisor of \(q^k-1\). The results below compute the
multiplicative order of \(q\) modulo \(S_k\), write down the positive cycle of
the seed \(1\) with its minimal period, describe the residue walk of
\(1\bmod S_k\), locate the trap state, and recover the cycle from the return
map. Every statement is about \(\N\) at the raw clock: one division or one
affine step is one instruction.

\section{The repunit sum}

\begin{lemma}[Basic identities]\label{Lem:RepunitBasic}
For all \(q,k\ge0\) one has \(S_0=0\), \(S_{k+1}=S_k+q^k\), \(S_1=1\), and
\(qS_k+1=S_{k+1}\). For all \(j,k\ge0\) one has
\(S_{j+k}=S_j+q^jS_k\), and for \(k\ge1\) one has
\(S_k=q^{k-1}+S_{k-1}\).
\end{lemma}

\begin{proof}
All identities expand the defining sum. The recurrence \(S_{k+1}=S_k+q^k\)
splits the range \(\{0,\dots,k\}\) into \(\{0,\dots,k-1\}\) and \(\{k\}\).
The addition rule splits \(\{0,\dots,j+k-1\}\) into \(\{0,\dots,j-1\}\) and
\(\{j,\dots,j+k-1\}\), and the second block contributes
\(q^j(1+q+\cdots+q^{k-1})=q^jS_k\). The formula \(qS_k+1=S_{k+1}\) is the
case \(j=1\) of the addition rule together with \(S_1=1\), and
\(S_k=q^{k-1}+S_{k-1}\) is the recurrence written backwards.
\end{proof}

\begin{lemma}[Divisor identity]\label{Lem:DivisorIdentity}
For \(q\ge1\) and all \(k\ge0\) one has \((q-1)S_k=q^k-1\), and consequently
\(S_k\mid q^k-1\).
\end{lemma}

\begin{proof}
The difference telescopes,
\[
(q-1)S_k=\sum_{j=0}^{k-1}\left(q^{j+1}-q^j\right)=q^k-1,
\]
because consecutive powers cancel in the finite sum. Divisibility follows
with quotient \(q-1\).
\end{proof}

\begin{lemma}[Growth]\label{Lem:Growth}
For \(k\ge1\) one has \(1\le S_k\) and \(0<S_k\). For \(q\ge2\) one has
\(S_k<S_{k+1}\), the sequence \(k\mapsto S_k\) is strictly increasing,
\(S_k+1\le q^k\), and \(S_k<q^k\).
\end{lemma}

\begin{proof}
The first two claims are immediate from \(S_k=1+\cdots+q^{k-1}\) with
\(k\ge1\), which contains the term \(1\). For strict growth,
\(S_{k+1}=S_k+q^k>S_k\) because \(q^k>0\). Monotonicity follows by induction
on the difference of the indices. The bound \(S_k+1\le q^k\) is proved by
induction on \(k\): it is \(0+1\le1\) for \(k=0\), and
\[
S_{k+1}+1=S_k+q^k+1\le2q^k\le q^{k+1}
\]
for \(q\ge2\), using the induction hypothesis and \(2\le q\). The strict
inequality \(S_k<q^k\) follows from \(S_k+1\le q^k\).
\end{proof}

\begin{theorem}[Exact order]\label{Thm:Order}
Let \(q\ge2\). Then \(q^k\equiv1\pmod{S_k}\) for all \(k\ge0\), and for
\(0<j<k\) one has \(q^j\not\equiv1\pmod{S_k}\). Hence for \(k\ge1\) the
multiplicative order of \(q\) modulo the nonzero divisor \(S_k\) is exactly
\(k\).
\end{theorem}

\begin{proof}
The congruence is Lemma~\ref{Lem:DivisorIdentity}. Suppose now that
\(0<j<k\) and \(S_k\mid q^j-1\). By Lemma~\ref{Lem:Growth} the sum \(S_j\)
satisfies \(S_j<q^j\) and \(S_j<S_k\), and by the addition rule
\(S_k=S_j+q^jS_{k-j}\) with \(k-j\ge1\), so \(q^j\le q^jS_{k-j}=S_k-S_j<S_k\).
Lemma~\ref{Lem:DivisorIdentity} gives \(q^j-1=(q-1)S_j\), hence
\(0\le q^j-1<S_k\). A positive multiple of \(S_k\) in that interval is zero,
so \(q^j=1\), which for \(q\ge2\) and \(j>0\) is impossible. Therefore
\(S_k\nmid q^j-1\) and the order is \(k\).
\end{proof}

\section{The unit cycle}

\begin{lemma}[Branches and the trap]\label{Lem:StepBranch}
For all \(q,D,n\) one has \(T_{q,D}(n)=n/D\) when \(D\mid n\) and
\(T_{q,D}(n)=qn+1\) when \(D\nmid n\); the value of an iterate after \(k+1\)
steps is the value after \(k\) steps followed by one more step; and
\(T_{q,D}(0)=0\) with \(T_{q,D}^{\,j}(0)=0\) for every \(j\).
\end{lemma}

\begin{proof}
The first two claims are the definition and the recursion of the iterate. For
the last one, \(D\mid0\) with quotient \(0\), so \(0\) is a fixed point, and
an iterate of a fixed point is that fixed point.
\end{proof}

\begin{theorem}[Trap uniqueness]\label{Thm:Trap}
If \(T_{q,D}(n)=0\), then \(n=0\).
\end{theorem}

\begin{proof}
If \(D\mid n\), then \(n/D=0\) forces \(n=0\). If \(D\nmid n\), then
\(qn+1=0\) is impossible because \(qn+1\ge1\).
\end{proof}

\begin{theorem}[The unit cycle]\label{Thm:Cycle}
Let \(q\ge2\) and \(k\ge1\), and put \(D=S_k\). For \(0\le j<k\) one has
\(T_{q,D}^{\,j}(1)=S_{j+1}\), and \(T_{q,D}^{\,k}(1)=1\). For
\(0<j<k\) one has \(T_{q,D}^{\,j}(1)\neq1\). Hence \(1\) lies on a positive
cycle of minimal period \(k\), whose states are \(S_2,\dots,S_k\) in that
order.
\end{theorem}

\begin{proof}
Induct on \(j<k\). For \(j=0\) the claim is \(S_1=1\), which is
Lemma~\ref{Lem:RepunitBasic}. For the step, let \(j+1\le k\) and suppose that
\(T^{\,j}(1)=S_{j+1}\). If \(j+1=k\), then the state is \(D\) itself, a
multiple of \(D\), so the division branch gives \(T^{\,k}(1)=D/D=1\). If
\(j+1<k\), then \(0<S_{j+1}<S_k=D\) by Lemma~\ref{Lem:Growth}, so the state is
not a multiple of \(D\) and the affine branch applies:
\[
T^{\,j+1}(1)=qS_{j+1}+1=S_{j+2}
\]
by Lemma~\ref{Lem:RepunitBasic}. This proves both the orbit formula and the
return. For minimality, let \(0<j<k\). By the formula just proved the state is
\(T^{\,j}(1)=S_{j+1}\) with \(2\le j+1\le k\), and Lemma~\ref{Lem:Growth}
gives \(S_{j+1}\ge S_2=1+q>1\), so \(1\) does not return earlier. All these
states are distinct by the strict monotonicity of Lemma~\ref{Lem:Growth},
which also shows that \(S_2,\dots,S_k\) are pairwise different and different
from \(S_1=1\).
\end{proof}

\section{The return map}

\begin{theorem}[Return map]\label{Thm:Return}
Let \(q\ge1\) and \(k\ge2\), and define \(R(n)=(q^{\,k-1}n+S_{k-1})/S_k\).
After \(k-1\) affine events from the seed \(1\) the state of the slice is
\(S_k\), and \(R(1)=1\).
\end{theorem}

\begin{proof}
By Theorem~\ref{Thm:Cycle} with \(j=k-1\) the state after \(k-1\)
instructions is \(S_k\), which is the divisor itself, so the following
instruction is a division and returns \(S_k/S_k=1\). The closed form is the
same computation: \(q^{\,k-1}+S_{k-1}=S_k\) by
Lemma~\ref{Lem:RepunitBasic}, hence
\[
R(1)=\frac{q^{\,k-1}\cdot1+S_{k-1}}{S_k}=\frac{S_k}{S_k}=1.
\]
The denominator is positive because \(S_k\ge1\) by Lemma~\ref{Lem:Growth}.
\end{proof}

\section{The residue walk}

\begin{theorem}[Residue walk]\label{Thm:ResidueWalk}
Let \(q\ge2\), let \(D=S_k\) and let \(\rho(r)=(qr+1)\bmod D\). For every
\(j\) one has \(\rho^{\,j}(0)=S_j\bmod D\); the walk is periodic,
\(\rho^{\,j+k}(0)=\rho^{\,j}(0)\) for every \(j\); and the residues
\(S_0,S_1,\dots,S_{k-1}\) modulo \(D\) are pairwise distinct.
\end{theorem}

\begin{proof}
The formula is an induction on \(j\): it holds at \(j=0\) because
\(S_0=0\), and the step is \(\rho(S_j)=(qS_j+1)\bmod D=S_{j+1}\bmod D\) by
Lemma~\ref{Lem:RepunitBasic}. For the period,
\[
\rho^{\,j+k}(0)=S_{j+k}\bmod D=\left(S_j+q^jS_k\right)\bmod D=S_j\bmod D
=\rho^{\,j}(0),
\]
using the addition rule and \(S_k=D\). For distinctness, let \(j,j'<k\) with
\(S_j\equiv S_{j'}\pmod D\). Lemma~\ref{Lem:Growth} gives \(S_j<D\) and
\(S_{j'}<D\), so both residues are already reduced and \(S_j=S_{j'}\);
strict monotonicity then gives \(j=j'\). In particular the accessible
residues are \(0,S_1,\dots,S_{k-1}\), and no other residue is reached from
\(0\).
\end{proof}

\section{Controls and scope}

\begin{remark}[Boundaries of the family]\label{Rem:Controls}
The hypotheses are needed. For \(q=1\) the divisor identity reads
\(0=0\) and \(1^{\,j}\equiv1\pmod{k}\) for every \(j\), so the order claim
fails and the rows using it assume \(q\ge2\). For \(k=1\) the divisor is
\(S_1=1\), every state is divisible by it, the map is the identity and the
cycle has length one; the return theorem assumes \(k\ge2\). For a divisor that
is not a repunit sum, for instance \(q=2\), \(D=5\), the raw orbit of the seed
\(1\) does not return to \(1\) within the tested horizon, while the residue
walk from \(0\) does close, because \(0\mapsto1\mapsto3\mapsto2\mapsto0\) is a
cycle of the residue map. The classical instance
\(q=3\), \(k=2\) gives \(D=4\) and the cycle \(1\mapsto4\mapsto1\).
\end{remark}

\begin{remark}[Scope]\label{Rem:Scope}
The statements classify the slice through the order of \(q\) modulo \(S_k\)
and the listed cycle. They do not claim that seeds outside the cycle diverge,
they say nothing about the density of the captured set, and they do not treat
general divisors \(D\) that are not repunit sums. The bounded Go certificate
of \texttt{go/RepunitDivisor.go} is an agreement check for the statements
above, not a replacement for them.
\end{remark}

\section{Formal correspondence}

\begin{table}[htbp]
\centering
\begin{tabular}{ll}
\toprule
Statement & Lean / controls \\
\midrule
Definition of \(S_k\), \(T\), \(\rho\), \(R\) &
\texttt{RepunitSum}, \texttt{RepunitStep}, \texttt{RepunitResidue}, \texttt{RepunitReturn} \\
Lemma~\ref{Lem:RepunitBasic} & \texttt{RepunitSumZero}, \texttt{RepunitSumSucc}, \texttt{RepunitSumOne}, \texttt{RepunitSumSuccMul}, \texttt{RepunitSumAdd}, \texttt{RepunitSumEqPowAdd} \\
Lemma~\ref{Lem:DivisorIdentity} & \texttt{RepunitSumMulPred}, \texttt{RepunitSumDvdPred} \\
Theorem~\ref{Thm:Order} & \texttt{RepunitSumModeq}, \texttt{RepunitSumModeqOrder} \\
Lemma~\ref{Lem:Growth} & \texttt{RepunitSumOneLe}, \texttt{RepunitSumPos}, \texttt{RepunitSumLtSucc}, \texttt{RepunitSumStrictMono}, \texttt{RepunitSumSuccLePow}, \texttt{RepunitSumLtPow} \\
Lemma~\ref{Lem:StepBranch} & \texttt{RepunitStepOfDvd}, \texttt{RepunitStepOfNotDvd}, \texttt{RepunitIterateSuccLast}, \texttt{RepunitStepZero}, \texttt{RepunitStepIterateZero} \\
Theorem~\ref{Thm:Trap} & \texttt{RepunitStepEqZero} \\
Theorem~\ref{Thm:Cycle} & \texttt{RepunitStepIterateOne}, \texttt{RepunitStepIterateOneCycle}, \texttt{RepunitStepIterateOneNe} \\
Theorem~\ref{Thm:Return} & \texttt{RepunitReturnFixesOne} \\
Theorem~\ref{Thm:ResidueWalk} & \texttt{RepunitResidueIterateZero}, \texttt{RepunitResidueIteratePeriod}, \texttt{RepunitResidueIterateInjective} \\
Remark~\ref{Rem:Controls} & control slice of \texttt{go/RepunitDivisor.go}; boundary statements above \\
Remark~\ref{Rem:Scope} & scope remarks for the statements above \\
\bottomrule
\end{tabular}
\end{table}

\end{document}
